%-------------------------------------------------------------------------------
\section{Rule Extraction and Lintability Determination}
\label{sec:methodology}
%-------------------------------------------------------------------------------

This section presents DLEP, our framework for Web PKI normative rule extraction and lintability analysis.
We describe the overall system architecture and the methods used in the key components.

\subsection{Architecture Overview}
\label{sec:method-overview}

Web PKI normative sources serve as the input to the framework. Through keyword recall and constrained extraction, candidate rules are converted into a structured representation. As shown in Fig.~\ref{fig:architecture}, the framework contains four modules.

(1)~Knowledge graph construction. To handle scattered definitions, inherited requirements, and cross-document references, this module organizes the collected sources into a traceable source context for rule extraction.

(2)~Deterministic context retrieval. To prevent the LLM from inventing unsupported context, this module supplies source-grounded context for each target section in a reproducible way.

(3)~DLEP. To balance completeness and precision, this module first recalls candidate normative rules; for precise extraction, the second layer uses an LLM to convert them into a structured representation.

(4)~Four-condition lintability framework. To determine whether a rule can be implemented as a static lint check, we design a four-condition lintability framework.

\begin{figure*}[tb]
  \centering
  \cicasfig{figures/fig01_architecture.png}{Architecture of the Dual-Layer
    Extraction Pipeline (DLEP), showing knowledge-graph construction, deterministic
    context retrieval, dual-layer rule extraction, and the four-condition
    lintability framework.}
  \caption{Architecture of Dual-Layer Extraction Pipeline (DLEP).}
  \label{fig:architecture}
\end{figure*}

\subsection{PKI Knowledge Graph Builder}
\label{sec:method-kg-builder}

The PKI knowledge graph builder is the offline component that converts heterogeneous Web PKI normative sources into the retrieval substrate used by DLEP.
It normalizes documents into a hierarchical structure with stable identifiers for specifications, sections, headings, and text units.
It then extracts explicit references such as section citations, RFC references, and policy cross-links, and links them to target nodes when the targets can be resolved unambiguously.
In parallel, it records certificate-field concepts and aliases, namely subjectAltName, dNSName, basicConstraints.cA, and policy OIDs, to map natural-language mentions to canonical certificate paths during later extraction.

PKI specifications form a densely interconnected cross-document network.
For example, RFC~5280 references multiple RFCs, while CABF~BR directly depends on RFC~5280 field definitions.
We model this knowledge as a multi-edge directed graph with eight node types, namely Specification, Section, Definition, CertificateField, Rule, Operation, Value, and Concept, and four source-backed relations, namely CONTAINS, DEFINES, REFERENCES, and APPLIES\_TO. Appendix~\ref{app:kg} summarizes these relations and their deterministic retrieval constraints.

This builder is deliberately conservative.
It stores only relations directly supported by document structure or textual evidence, while leaving normative judgments such as conflict, override, and lintability to later deterministic stages.
Normative-judgment relations such as CONFLICTS\_WITH and OVERRIDES are produced independently by the rule engine and are not included in the retrieval graph.
The resulting graph serves as a traceable context index rather than as an oracle. It helps the extractor identify definitions, inherited constraints, and referenced rules, but does not determine whether a requirement is lintable.

\subsection{Deterministic Context Retrieval}
\label{sec:method-graphrag}

To address cross-document references, DLEP organizes context deterministically rather than relying on the LLM to infer relations.
We implement a GraphRAG~\cite{edge2024local} retrieval module that treats the knowledge graph of \S\ref{sec:method-kg-builder} as a fixed input and retrieves context through bounded subgraph queries.

Given a target Section node, the retrieval module performs a BFS expansion over the $k$-hop neighborhood using the source-backed CONTAINS, DEFINES, REFERENCES, and APPLIES\_TO relations.
It then assembles LLM context by relation-type priority and provides term definitions, field metadata, and referenced sections. The retrieval process uses neither vector similarity nor LLM inference, and produces no independent judgment.

\subsection{DLEP}
\label{sec:method-dual-layer}

\begin{table}[tb]
\centering
\caption{RFC~2119 keyword tiers.}
\label{tab:rfc2119}
\setlength{\tabcolsep}{4pt}
\begin{tabular}{@{}L{0.27\linewidth} L{0.30\linewidth} L{0.33\linewidth}@{}}
\hline
Tier & Keywords & Semantics \\
\hline
Mandatory      & MUST, SHALL, REQUIRED   & Absolute requirement; a violation is non-compliance \\
Prohibition    & MUST NOT, SHALL NOT     & Absolute prohibition \\
Recommendation & SHOULD, SHOULD NOT, RECOMMENDED & Strongly recommended or discouraged; justified deviation permitted \\
Optional       & MAY, OPTIONAL           & Fully optional \\
\hline
\end{tabular}
\end{table}

\begin{table}[tb]
\centering
\caption{IR core four-tuple fields.}
\label{tab:ir-core}
\setlength{\tabcolsep}{4pt}
\footnotesize
\begin{tabular}{@{}L{0.18\linewidth} L{0.74\linewidth}@{}}
\hline
Field & Meaning \\
\hline
subject    & Certificate field path (e.g., extensions.\allowbreak subjectAltName.\allowbreak dNSName), resolved by FieldResolver \\
obligation & RFC~2119 keyword (MUST/SHALL/SHOULD/\dots) \\
predicate  & Assertion type (must\_be\_present, conform\_to, equal, \dots) \\
constraint & constraint value or pattern \\
\hline
\end{tabular}
\end{table}

To address the LLM controllability challenge, we design a dual-layer extraction pipeline (DLEP) that preserves the LLM's semantic understanding while keeping the extraction process auditable.

Layer~1 performs deterministic recall.
Web PKI specifications commonly express normative requirements using the RFC~2119 keywords listed in Table~\ref{tab:rfc2119}.
RFC~8174 further specifies that these keywords carry normative force only when capitalized, enabling deterministic keyword matching.
Since ETSI~EN~319 does not follow this convention, its candidates are recalled through lowercase matching.

To maximize recall, Layer~1 operates in three passes.
Pass~1 directly matches RFC~2119 keywords.
Pass~2 detects nested structures so subordinate rules inherit the parent obligation level.
Pass~3 captures normative statements that do not follow the standard RFC~2119 form.

Layer~2 performs controlled semantic parsing using an LLM, but restricts the model to a constrained language-understanding role.
Instead of generating lint rules or compliance judgments directly, Layer~2 converts normative text into a JSON-serialized intermediate representation (IR), while all normative reasoning is deferred to later deterministic stages.

The IR decouples semantic extraction from downstream tasks, allowing the same extraction result to support multiple uses, as shown in Fig.~\ref{fig:ir-downstream}.
Each IR record is serialized as a JSON object with 33 fields.
Its core component, shown in Table~\ref{tab:ir-core}, is the four-tuple
$\mathrm{IR} = \langle\text{subject}, \text{obligation}, \text{predicate}, \text{constraint}\rangle$,
together with 29 extension fields summarized in Appendix~\ref{app:ir}.

This IR-based design has three advantages over direct end-to-end generation of lint checks from raw text.
First, citation, constraint, and obligation information are stored separately, improving traceability and conflict handling.
Second, the IR serves as deterministic input to the four-condition lintability analysis in \S\ref{sec:method-lintability}, making lintability decisions reproducible rather than model-dependent.
Third, once a rule is determined to be lintable, generating static lint checks from the IR is more stable and interpretable because key elements such as field paths, assertion types, and constraint values are already explicitly structured.

To constrain LLM behavior, Layer~2 enforces a predefined schema and category set.
Every extracted rule must map to explicit source text, and certificate field hierarchies must follow the ASN.1 path definitions in RFC~5280.
The IR generation process is also staged. The model first determines the rule category and then fills in the remaining fields, which reduces errors compared with generating all fields in a single step.
Appendix~\ref{app:ir} provides the category definitions, correction rules, and tuning details.

\begin{figure}[tb]
  \centering
  \cicasfig{figures/fig02_ir_downstream.png}{IR-driven ``extract once, apply
    many'' architecture.}
  \caption{IR-driven ``extract once, apply many'' architecture.}
  \label{fig:ir-downstream}
\end{figure}

\subsection{Normative Rule Lintability Decision}
\label{sec:method-lintability}
\label{sec:lintability}

\begin{table*}[tb]
\centering
\caption{The four lintability conditions, each a Boolean function of one IR field.}
\label{tab:four-conditions}
\footnotesize
\setlength{\tabcolsep}{4pt}
\renewcommand{\arraystretch}{1.2}
\begin{tabular}{@{}l l L{0.36\linewidth} L{0.30\linewidth}@{}}
\hline
ID & IR field & Value domain & Condition \\
\hline
$C_1$      & obligation         & MUST, SHALL, REQUIRED, SHOULD, RECOMMENDED, MUST NOT, SHALL NOT, SHOULD NOT, MAY, OPTIONAL & obligation $\notin$ \{MAY, OPTIONAL\} \\
$C_2$      & assertion\_subject & Certificate, CA, CrossArtifact & assertion\_subject $=$ Certificate \\
$C_3$      & enforcement\_phase & Encoding, Validation, Comparison, Processing & enforcement\_phase $=$ Encoding \\
$\neg C_4$ & rule\_category     & encoding\_constraint, comparison, definition, capability, algorithm\_ref, display, \dots & rule\_category $\notin$ \{definition, capability, algorithm\_ref, display, \dots\} \\
\hline
\end{tabular}
\end{table*}

Before generating code, we decide the lintability of each normative rule, namely whether it admits such a static lint check. Rather than asking the LLM to make this judgment directly, we compute it from the IR that the LLM has already produced. Three observations motivate this design. First, the same four pieces of information determine the lintability of normative rules in practice, namely the deontic strength of the obligation, who is obligated, when the obligation applies, and what type of rule it is. Second, each of the four can be encoded as a single IR field with a small, closed value set, so the lintability decision becomes a Boolean function of four discrete fields. Third, separating the four conditions makes any wrong label traceable to exactly one IR field rather than to an opaque end-to-end LLM call.

Table~\ref{tab:four-conditions} lists the four IR fields, their value domains, and the conditions they decide. The deontic level obligation comes from the IR core four-tuple in Table~\ref{tab:ir-core}. The remaining three fields, assertion\_\allowbreak subject, enforcement\_\allowbreak phase, and rule\_\allowbreak category, are key extension fields. From the enumeration of rule\_\allowbreak category, we designate
\begin{gather*}
\mathcal{N} = \{definition,\,capability,\\
algorithm\_ref,\,display,\,\dots\}
\end{gather*}
as the subset whose rule type carries no static certificate lint check on the certificate. This subset covers term definitions, CA capability statements, delegations to external algorithm specifications, UI rendering rules, and similar non-encoding rule types.

The four conditions in Table~\ref{tab:four-conditions} are necessary for lintability, and each is a Boolean function of exactly one IR field:
\begin{equation}\label{eq:lintable}
\mathrm{lintable}(r) \;\Longrightarrow\; C_1(r)\wedge C_2(r)\wedge C_3(r)\wedge \neg C_4(r),
\end{equation}
\begin{align*}
C_1(r) &\equiv \mathrm{is\_normative}(r.\text{obligation}), \\
C_2(r) &\equiv r.\text{assertion\_subject} = Certificate, \\
C_3(r) &\equiv r.\text{enforcement\_phase} = Encoding, \\
C_4(r) &\equiv r.\text{rule\_category} \in \mathcal{N}.
\end{align*}
The obligation field also fixes the lint severity. MUST-class obligations map to Error, and SHOULD-class obligations map to Warning. Severity is therefore not a second classifier, but a direct read of obligation.

The three non-deontic conditions encode three orthogonal boundaries. $C_2$ is the subject boundary, separating obligations on the certificate from obligations on the CA, relying party, or surrounding ecosystem. $C_3$ is the runtime boundary, separating obligations observable in the issued bytes from those that fire during chain processing, name comparison, revocation handling, or CAA retrieval. $\neg C_4$ is the process boundary, separating state assertions from terms in $\mathcal{N}$.

Since each $C_i$ is a deterministic function of one IR field, the lintability
decision---which adopts the conjunction $C_1\wedge C_2\wedge C_3\wedge\neg C_4$ of
these necessary conditions as its operational criterion---is computed from
the IR rather than predicted by the LLM. The same normative rule IR yields the same label bit-for-bit on every run, and any misclassification is traceable to a specific Layer~2 field assignment.
